% =============================================================================
%  Overlapping Temporal Decomposition (OTD) — IEEE Transactions style
%  Preamble:
%     \usepackage{amsmath,amssymb,bm}
%     \usepackage{algorithm}
%     \usepackage[noend]{algpseudocode}
% =============================================================================

\section{Overlapping Temporal Decomposition}
\label{sec:otd}

The multi-period OPF is coupled across time only through the scalar
state-of-charge (SOC) chain
\begin{equation}
\label{eq:soc_chain}
B_{t,j} = B_{t-1,j} + \eta_{c,j}P^{c}_{t,j} - \eta_{d,j}^{-1}P^{d}_{t,j},
\;\; t\!\in\!\mathcal{T},\, j\!\in\!\mathcal{B},
\end{equation}
with horizon $\mathcal{T}\!=\!\{1,\dots,T\}$ and battery set $\mathcal{B}$.
We exploit this thin temporal coupling through an \emph{overlapping
temporal decomposition} (OTD) of Schwarz type: the horizon is partitioned
into $P$ contiguous cores, each core is enlarged on \emph{both} sides by
an overlap region, and the resulting subproblems are solved in parallel
while exchanging only the boundary SOC values and the chain duals.

% -----------------------------------------------------------------------------
\subsection{Window Construction}
\label{sec:otd_windows}

Let $b\!=\!\lfloor T/P\rfloor$ and let $\Delta\!\ge\!1$ be a user-chosen
overlap. Window $i\!\in\!\{1,\dots,P\}$ is described by a core interval
$[c^{s}_{i},c^{e}_{i}]$ and a solve (extended) interval
$[w^{s}_{i},w^{e}_{i}]$:
\begin{equation}
\label{eq:windows}
\begin{aligned}
c^{s}_{i} &= (i\!-\!1)b\!+\!1, & c^{e}_{i} &= ib,\\
w^{s}_{i} &= \max(c^{s}_{i}\!-\!\Delta,1), & w^{e}_{i} &= \min(c^{e}_{i}\!+\!\Delta,T).
\end{aligned}
\end{equation}
The cores tile $\mathcal{T}$ without gaps, while the extended intervals
overlap their neighbours by up to $\Delta$ steps on each side. A larger
$\Delta$ improves boundary contraction and tightens the global optimality
gap, at a proportional increase in subproblem size; $\Delta\!=\!b$ is a
robust default.

% -----------------------------------------------------------------------------
\subsection{Boundary Variables, Exchange, and Termination}
\label{sec:otd_bc}

In order to solve any local subproblem, the values at its temporal
boundaries must be pre-specified: at iteration $k\!=\!0$ they are supplied
by an initialisation, and at $k\!\ge\!1$ they are taken from the most
recent solutions of the neighbouring windows. Two boundaries are
\emph{fixed} throughout the Schwarz loop and are not exchanged:
\begin{equation}
\label{eq:fixed_bc}
\widetilde{B}^{\mathrm{L}}_{1,j}\;\equiv\;B_{0,j},
\qquad
\widetilde{B}^{\mathrm{R}}_{P,j}\;\equiv\;B_{0,j},
\end{equation}
namely the left boundary of window $i\!=\!1$ (the given initial SOC) and
the right boundary of window $i\!=\!P$ (the cyclic terminal SOC, set equal
to the initial SOC of window $1$). All remaining interface variables are
initialised to $B_{0,j}$ (and the dual hint to zero) and refreshed at
every iteration as described below.

\subsubsection*{Per-window information flow}
Consider an interior window $i\!\in\!\{2,\dots,P\!-\!1\}$ at iteration
$k$. Its solve on $[w^{s}_{i},w^{e}_{i}]$ requires three boundary
parameters from its neighbours and, in turn, produces three quantities
that its neighbours will consume.

\emph{Parameters received from the upstream window $i\!-\!1$:}
\begin{equation}
\label{eq:left_recv}
\widetilde{B}^{\mathrm{L},(k)}_{i,j}
\;=\; B^{(i-1),(k)}_{w^{s}_{i},\,j},
\end{equation}
that is, the value of the upstream primal SOC trajectory evaluated at the
time step $t\!=\!w^{s}_{i}$. This time step lies inside the right overlap
of window $i\!-\!1$, so its value is well-defined.

\emph{Parameters received from the downstream window $i\!+\!1$:}
\begin{align}
\widetilde{B}^{\mathrm{R},(k)}_{i,j}
&= B^{(i+1),(k)}_{w^{e}_{i},\,j},
\label{eq:right_recv_B}\\
\lambda^{\mathrm{R},(k)}_{i,j}
&= \mu^{(i+1),(k)}_{w^{e}_{i}+1,\,j},
\label{eq:right_recv_lam}
\end{align}
i.e., the downstream primal SOC at the right cut $t\!=\!w^{e}_{i}$ and the
Lagrange multiplier of the SOC chain~\eqref{eq:soc_chain} in window
$i\!+\!1$ at $t\!=\!w^{e}_{i}\!+\!1$. The latter is the marginal cost,
according to the downstream subproblem, of one additional unit of SOC
delivered at the right cut. Injecting~\eqref{eq:right_recv_lam} into the
objective of window $i$ therefore allows the local solve to ``see'' the
future cost of its own choice of terminal SOC, mimicking the price
information that the centralised problem would have propagated.

\emph{Quantities produced and shared after the local solve at $k\!+\!1$:}
\begin{equation}
\label{eq:produce}
\underbrace{B^{(i),(k+1)}_{w^{s}_{i+1},\,j}}_{\text{to }i+1},
\quad
\underbrace{B^{(i),(k+1)}_{w^{e}_{i-1},\,j}}_{\text{to }i-1},
\quad
\underbrace{\mu^{(i),(k+1)}_{w^{e}_{i-1}+1,\,j}}_{\text{to }i-1}.
\end{equation}
The two boundary windows specialise this exchange: window $1$ receives no
left-side data (its left BC is the constant $B_{0,j}$) and only emits
data to window $2$; symmetrically, window $P$ emits data to window
$P\!-\!1$ but receives no right-side data.

\subsubsection*{Boundary update}
Let $\widehat{B}^{\bullet,(k+1)}_{i,j}$ and
$\widehat{\lambda}^{\mathrm{R},(k+1)}_{i,j}$ ($\bullet\!\in\!\{\mathrm{L},\mathrm{R}\}$)
denote the fresh candidate values assembled
from~\eqref{eq:left_recv}--\eqref{eq:produce} after all $P$ subproblems
return. The \emph{default} update is full assignment,
\begin{equation}
\label{eq:full_update}
\widetilde{B}^{\bullet,(k+1)}_{i,j} \leftarrow
\widehat{B}^{\bullet,(k+1)}_{i,j},
\qquad
\lambda^{\mathrm{R},(k+1)}_{i,j} \leftarrow
\widehat{\lambda}^{\mathrm{R},(k+1)}_{i,j},
\end{equation}
i.e.\ each window adopts the new boundary values supplied by its
neighbours. Full assignment is observed to converge on all
well-conditioned cases. In the rare event that the primal SOC boundary
\emph{oscillates} between consecutive iterations — a known Schwarz
failure mode when the local problems exhibit flat directions at the
cut — a relaxation factor $\alpha\!\in\!(0,1]$ is invoked to dampen the
update by blending the new and previous values:
\begin{equation}
\label{eq:relax}
\widetilde{B}^{\bullet,(k+1)}_{i,j} =
\alpha\widehat{B}^{\bullet,(k+1)}_{i,j}
+ (1\!-\!\alpha)\widetilde{B}^{\bullet,(k)}_{i,j}.
\end{equation}
The dual hint is always assigned directly, since it acts as a
terminal-cost predictor rather than a coupled primal iterate. Full
assignment is recovered for $\alpha\!=\!1$; values $\alpha\!\in\![0.5,0.8]$
are typical when damping is required.

\subsubsection*{Termination criterion}
The Schwarz loop terminates when the maximum mismatch between the
freshly computed candidates and the previous iterate, taken over all
interfaces, batteries, and sides, falls below a tolerance
$\varepsilon\!>\!0$:
\begin{equation}
\label{eq:term}
\Delta B^{(k)} \,:=\, \max_{i,j,\,\bullet}
\big|\widehat{B}^{\bullet,(k+1)}_{i,j} - \widetilde{B}^{\bullet,(k)}_{i,j}\big|
\;<\; \varepsilon,
\end{equation}
with an iteration cap $K_{\max}$ as safeguard. Upon termination, the
global solution is stitched from the \emph{core} portions of the
converged windows,
\begin{equation}
\label{eq:stitch}
\bm{x}^{\star}_{t} = \bm{x}^{(i)}_{t},
\quad t\!\in\![c^{s}_{i},c^{e}_{i}],\; i=1,\dots,P,
\end{equation}
so that the overlap timesteps are discarded once they have served their
information-passing role. Note that the per-iteration message volume
is $\mathcal{O}(P|\mathcal{B}|)$, independent of $T$ and of the network
size.

% -----------------------------------------------------------------------------
\subsection{Local Subproblem}
\label{sec:otd_obj}

Let $\bm{x}_{i}$ collect the local OPF variables of window $i$ on
$[w^{s}_{i},w^{e}_{i}]$, and let $\mathcal{X}_{i}$ denote its local
feasibility set (power balance, branch and voltage limits, SOCP
relaxation, battery limits, and the SOC chain~\eqref{eq:soc_chain}
anchored on the left by~\eqref{eq:left_recv} or, for $i\!=\!1$,
by~\eqref{eq:fixed_bc}). The local economic cost reads
\begin{equation}
\label{eq:local_cost}
f_{i}(\bm{x}_{i}) = \!\!\sum_{t=w^{s}_{i}}^{w^{e}_{i}}\!\!
   c_{t}\!\!\!\sum_{(b,\phi)\in\mathcal{S}}\!\!\! P^{\mathrm{subs}}_{t,b,\phi}
 \;+\; \alpha_{\mathrm{scd}}\!\!\sum_{t,j}\!
   \big(a^{c}_{j}P^{c}_{t,j}+a^{d}_{j}P^{d}_{t,j}\big),
\end{equation}
with $c_{t}$ the energy price, $\mathcal{S}$ the substation phase set,
$a^{c}_{j}\!=\!1\!-\!\eta_{c,j}$, $a^{d}_{j}\!=\!\eta_{d,j}^{-1}\!-\!1$,
and $\alpha_{\mathrm{scd}}$ a small simultaneous charge--discharge
regulariser. At Schwarz iteration $k$ window $i$ solves
\begin{equation}
\label{eq:otd_sub}
\bm{x}_{i}^{(k+1)} \in
\arg\!\min_{\bm{x}_{i}\in\mathcal{X}_{i}}\,
f_{i}(\bm{x}_{i}) + \phi^{\mathrm{R}}_{i,k}(\bm{x}_{i}),
\end{equation}
with the right-boundary augmented-Lagrangian term
\begin{equation}
\label{eq:al_term}
\phi^{\mathrm{R}}_{i,k} \!=\! \sum_{j\in\mathcal{B}}\!
\Big[\lambda^{\mathrm{R},(k)}_{i,j} B^{(i)}_{w^{e}_{i},j}
\!+\!\rho\big(B^{(i)}_{w^{e}_{i},j}\!-\!\widetilde{B}^{\mathrm{R},(k)}_{i,j}\big)^{\!2}\Big],
\end{equation}
which is suppressed for the terminal window ($i\!=\!P$, whose right SOC
is hard-fixed by~\eqref{eq:fixed_bc}). The linear term embeds the
downstream marginal price of energy into the local objective, while the
quadratic term enforces consensus on the boundary SOC. The penalty
$\rho\!>\!0$ is the single tuning parameter; $\rho\!=\!\mathcal{O}(1)$
yields monotone contraction in our experiments, whereas large $\rho$
over-weights the consensus term and induces oscillation of the AL
multiplier.

\emph{Remark (ISOCP).} When the SOCP relaxation is tightened iteratively
(ISOCP), refinement cuts are \emph{not} added inside the Schwarz loop,
as they perturb the duals $\mu^{(i)}_{t,j}$ between iterations and break
stationarity. A single ISOCP refinement pass is performed per window
after the boundaries have converged.

% -----------------------------------------------------------------------------
\subsection{Parallel Algorithm}
\label{sec:otd_algo}

The complete procedure is given in Algorithm~\ref{alg:otd}. Each window is
assigned a dedicated worker process that builds its local OPF model once
at start-up and thereafter receives only the small boundary tuple at
each iteration; the per-iteration communication cost is therefore
negligible relative to the local solve.

\begin{algorithm}[t]
\caption{Parallel OTD--Schwarz}
\label{alg:otd}
\begin{algorithmic}[1]
\Require horizon $T$, partitions $P$, overlap $\Delta$, penalty $\rho$,
         tolerance $\varepsilon$, cap $K_{\max}$, initial SOC $B_{0}$.
\Ensure global solution $\bm{x}^{\star}$.

\vspace{2pt}\Statex \textit{Setup.}
\State Build windows $\{(c^{s}_{i},c^{e}_{i},w^{s}_{i},w^{e}_{i})\}_{i=1}^{P}$
       from~\eqref{eq:windows}.
\State Fix $\widetilde{B}^{\mathrm{L}}_{1}\!=\!\widetilde{B}^{\mathrm{R}}_{P}\!=\!B_{0}$;
       initialise all other $\widetilde{B}^{\mathrm{L}}_{i},
       \widetilde{B}^{\mathrm{R}}_{i}\!\leftarrow\!B_{0}$ and
       $\lambda^{\mathrm{R}}_{i}\!\leftarrow\!\bm{0}$.
\State \textbf{In parallel}, spawn one worker per window and build the local
       model~\eqref{eq:otd_sub} from $\mathcal{D}_{i}$.

\vspace{2pt}\Statex \textit{Schwarz loop.}
\For{$k=1,\dots,K_{\max}$}
  \State \textbf{Send} the current boundary tuple
         $(\widetilde{B}^{\mathrm{L}}_{i},\widetilde{B}^{\mathrm{R}}_{i},
         \lambda^{\mathrm{R}}_{i})$ to every worker.
  \State \textbf{In parallel}, each worker solves~\eqref{eq:otd_sub}
         and returns its primal SOC $B^{(i)}_{\cdot,j}$ and chain dual
         $\mu^{(i)}_{\cdot,j}$.
  \State \textbf{Assemble} candidate boundaries
         $\widehat{B}^{\mathrm{L}}_{i},\widehat{B}^{\mathrm{R}}_{i},
         \widehat{\lambda}^{\mathrm{R}}_{i}$
         from~\eqref{eq:left_recv}--\eqref{eq:right_recv_lam}.
  \State \textbf{Check} $\Delta B^{(k)}$ via~\eqref{eq:term};\;
         \textbf{if} $\Delta B^{(k)}<\varepsilon$ \textbf{then break}.
  \State \textbf{Update} interface variables: full
         assignment~\eqref{eq:full_update} by default,\;
         relaxation~\eqref{eq:relax} with $\alpha\!<\!1$ if oscillation
         is detected.
\EndFor

\vspace{2pt}\Statex \textit{Refinement and output.}
\If{ISOCP enabled}
  \State \textbf{In parallel}, each worker runs one ISOCP refinement at
         the frozen boundaries.
\EndIf
\State Stitch $\bm{x}^{\star}$ from the core slices via~\eqref{eq:stitch}.
\State Shut down workers and \Return $\bm{x}^{\star}$.
\end{algorithmic}
\end{algorithm}

\noindent The number of Schwarz iterations is essentially independent of
$T$ and grows mildly with $P$ for $\Delta$ chosen comparable to $b$. The
wall-clock speed-up over the centralised multi-period OPF therefore
approaches $P$ as $T$ grows, bounded only by the one-shot model-build
overhead and the boundary-exchange latency.
